\chapter[Sermon 75. The Godly Vision Is Yet More Plainly Declared.]{The Godly Vision Is Yet More Plainly Declared.}
\label{sermon:075}
\markright{Sermon 75. The Godly Vision Is Yet More Plainly Declared.}

The seven heads are seven mountains on which the woman sits; they are also seven kings. Five have fallen, one is present, and the other has not yet come. When he comes, he will continue for a time. And the beast that was, and is not, is also the eighth, and one of the seven, and will go into destruction. The ten horns which you saw are ten kings who have not yet received the kingdom, but will receive power as kings for one hour with the beast. These have one mind, and they will give their power and strength to the beast. These will fight against the Lamb, and the Lamb will overcome them. For he is Lord of all lords and King of all kings, and those who are on his side are called chosen and faithful.

The angelic interpreter proceeds to declare to John the mystery of the beast displayed, and of her judgment, in stages. And at this time, he explains three things: what the seven heads signify; why he said of the beast, “he was and is not”; and what the ten horns represent.

\index{John, Saint}\index{Rome}He explains two approaches. First, by seven mountains, upon which the woman sits, whom at the end of the chapter he calls the great City: namely, the great city, which is commonly called Rome, because it stands upon seven mountains. Moreover, the beast has seven heads, because it has had many times seven kings. Of this I spoke also in Chapter 13. At this point, he counts the seven kings, and there is no doubt that he is speaking of Rome. This, I suppose, is the Lord’s chief intent in these matters. For he could not speak more explicitly unless he had expressed the name of Rome itself; but the name of Babylon we heard expressed before. Five, he says, have fallen, namely, since the deadly wound was inflicted, in the death of Nero, within fourteen years. For immediately after Nero, Galba began to reign; who being slain, Otto reigned; which after killing himself, Vitellius succeeded, and he also was killed by the Flavians. For after him, Flavius Vespasian was emperor; after whom Titus, the best prince of all. And these five fell within fourteen years. He adds, and one of those is reigning now, namely, the sixth in order, Domitian, the son of Vespasian and brother to Titus, a most ungracious man, who persecuted the faithful and had condemned John the Apostle into exile. Another, says John, has not yet come: namely, Cocceius Nerva. For after he came to the empire, and lived most virtuously and righteously governed the empire, he did not remain long. For when he had reigned one year, three months, and nine days, he died. And thus much hitherto concerning the seven kings and the seven heads of the beast. These things, so certain, pertain not so much to the exposition of this passage as to the consolation of the faithful: who here may clearly perceive how empires consist in the hand and providence of God almighty, who knows his own and has care of the godly, although they may seem, by reason of their grievous persecutions and cruel torments, to be neglected by God.

Consequently, he explains why he said the beast was, and is not: namely, because he was the eighth Roman emperor, Ulpius Trajan. For he was the eighth from the empire wounded by Nero. Trajan was of the seventh, that is to say, was adopted by Nerva the seventh emperor. And until this point, the Roman Empire had been governed, first in deed by Caesars, then by the noblest citizens of Rome. But concerning this Trajan, who succeeded Nerva, the writers of histories say that he was the first foreigner to rule the empire, for he was a Spaniard. Therefore, the empire was or has been in the hands of the Romans, now it is so no more. For a Spaniard succeeds, so that the empire now seems as though it might be called Roman Spanish. And inasmuch as Trajan persecuted Christ and his members, he also went into perdition. And let no man think that this was the only and sole cause wherefore John said, concerning Trajan, that he was, and is not. For he has pronounced expressly, and he is the eighth, as though he should signify that there be other causes also, for the which it was said that the Roman Empire was, and now is not, whereof is spoken before.

\index{Apocalypse (Book of)}\index{Primasius}\index{Papacy}\index{Daniel (Book of)}Hereafter follows also the exposition of the ten horns. And the same horns are here recited, which are spoken of in the seventh of Daniel, and in the thirteenth of the Apocalypse. Neither is there any cause why one should superstitiously cling to the tenth number. For in the fourteenth of Numbers, the Lord says how he has been tempted ten times by the Israelites, for many times. Here is signified therefore, how the Roman Empire shall be dispersed into many kingdoms. For whether you say kings or kingdoms, the matter is all one. Doubtless the Roman Empire beginning to fall to decay, there sprang up kings in the East and West, which invaded the Roman Empire, Persians, Goths, Vandals, Lombards, and I know not what others; at the last in Spain, France, Hungary, I speak not of Africa and Asia, were found diverse kings, and the Roman monarchy ceased. Of these kings the Angel warns us for diverse causes. These, says he, have not yet received the kingdom. For whilst John wrote the Apocalypse, Domitian ruled, and the Roman Empire was yet mighty and strong, and so remained for certain ages. When therefore did they receive their kingdom? They receive, says he, power as kings at one hour with the beast, namely the second. For these things cannot be understood of the first and old Roman Empire. And Primasius, expounding this place, admonishes that an hour here is taken for a time present. Therefore at the same time, the beast, that new Empire grows up and increases, and the kings receive might and power. For the decay of the old Empire was the strength of kings, and of the new Papal Empire. And in deed the emperor Phocas commanded the church of Rome, and the Bishop thereof to be head of churches. Which gave a certain beginning to the Popes dominion, as also in the thirteenth chapter. I have recited: which he obtained at the length more fully under king Pippin and other Princes of France and Germany, but Nüchelus speaking of the Empire of Phocas in the twenty-first Generation, says the enemies of the Roman Empire, by the slothfulness and cowardice of Emperors, had taken away in the West country with the Islands of Germany, France, Spain, Hungary, Slavonia, and a good part of Italy, and thereto a great part of Africa; and in the East parts, Cacannus of Thrace, King of Huns invaded the Iberians, Armenians, Arabians, Dardanes, and the middle parts of Macedonie and Greece. And the Persians in a manner possessed all Assyria, the Saracens destroyed Egypt; fie for shame, our strength has so failed us through riot, covetousness, and voluptuousness, that the Roman Empire stood then only in name. Hitherto he. The same things have we discoursed more at large in the thirteenth chapter of this work. And verily Daniel shows how among those ten horns, one other little horn should grow up, which should strike off three, and take their place, and reign wantonly, cruelly, and wickedly. Wherefore the Papal Empire, and those sundry kingdoms grew up in a manner about one and the same time.

\index{Augustine, Saint}He also shows what kind of kingdoms those will be, and how they will behave toward that final beast, namely, the church of Rome: he says they all have one opinion; they believe one thing, and are of the same religion. He speaks chiefly of the western kings. For they all receive the decrees of the Bishop of Rome, and honor them, as most obedient children of the most sacred and holy church of Rome. They will deliver to the beast their power, their authority, or kingdom, for they submit themselves to the See of Rome. If the church of Rome has need of an army or force of arms, the kings send their power gladly to him, which the noble kingdom of Bohemia felt about a hundred years since, though it were to no great advantage and beautiful triumphs of the invaders. Yes, moreover, they acknowledge themselves to owe homage and fealty to the most holy and supreme Bishop in all the world. Hereunto chiefly pertains what Augustine, as recorded by Steuchus in his book against Laurence Valla, concerning the Donation of Constantine, in section 94, has written in this way: Gregory the Seventh to Géza, king of Hungary: “We suppose it is not unknown to you that the kingdom of Hungary, like as other most noble realms also, ought to be in the state of its own liberty, neither that it ought to be subject to any king of another realm, save to the holy and universal mother church of Rome, which has her subjects, not as servants, but as children.” Hereunto Steuchus adds: “Hear with what government the church rules, that she may sustain her subjects, not as servants, but as children. She does not displace kings from their possessions, but permits them to reign as her sons; while they are reigning, she reigns herself also.” Nevertheless, she will be known as Queen and Lady.

You hear how all the most noble realms are subject to the Apostolic See. Even there he shows that the most noble kingdoms of Spain, France, England, Denmark, Russia, Croatia, Dalmatia, Aragon, Sardinia, Portugal, Bohemia, Swabia, and Norway are subject and tributaries to the church of Rome. In section 97, he adds moreover: although the kings reigned and continued in possession, they are wont to acknowledge her as Queen, true Lady, and giver of their kingdoms. And in section 105, the old records of all Popes are full of high authority, whereby they have, with their empires, governed the whole world, having the rule and order of all lands, which power and authority that impudent praiser of the Roman See is not ashamed to call omnipotent or almighty.

\index{Antichrist}And undoubtedly we see today, great embassies sent to Rome by the newly elected and crowned western kings, intending to pay homage to the Pope, whom they call the head of Antichrist, and to offer what they term due obedience. Therefore, he called them before not kings absolutely, but as kings. For they acknowledge a superior, and are even as it were servants or wards of the servant of servants. Of whom he has made verses: “The common people brought from the far ends of the world, The servant of servants, O Rome, you are now their lord.”

\index{Justification}Hereunto the Apostle adds a thing yet more grievous. These kings, I mean the confederates of the Pope, and obedient children of the Church of Rome, endowed with the spirit of the beast, shall fight with the Lamb. Whereby is signified the tyranny which kings, and princes, and certain other states of the Roman Empire do practice, and long have practiced against Christ and his gospel. Concerning the Lamb we have already spoken enough before. John the Baptist, pointing with his finger to Christ, says: behold the Lamb of God, which takes away the sins of the world. Therefore shall the Roman princes fight, not against Christ himself, for they will be Christians, but against the Lamb, that is, the sanctification, justification and satisfaction of Christ. For if any man say at this day, that the son of God is most holy, by whom alone sins are forgiven, and we are sanctified: and say not also, that the Bishop of Rome is most holy also, which purges by pardons granted, but shall say rather, that pardons are plain deceitfulness, and the Pope most unclean of all: he shall doubtless neither be taken for a right Catholic, neither shall he be spared for confessing the Lamb of God. If any man shall confess that justification is only in the son of God alone, and that men are justified by faith only, and not also by our works and merits, he shall be carried to death or to prison, neither shall the confession of the Lamb of God avail him anything. If any man shall say, that he is fully purged through the only oblation of Christ on the cross, as of a lamb without spot, and sacrificed from the beginning, neither that he needs any popish Masses, whereby the priests boast that they make a daily offering for the sins of the quick and dead, which in deed is both false and blasphemous, he is straightways hurried to prison, and from thence drawn to the stake and burned. We cannot deny but that this is true, saying there are at this day innumerable examples of Romish kings and princes in this behalf. We shall not need therefore to fetch our exposition far off, how these kings, which wholly depend on the Pope, shall fight with the Lamb. I speak here nothing of others, which cleave whole unto Christ.

And therefore, for comfort, what follows is joined to this: the Lamb will overcome them. For although Papist kings and princes seem to overcome the saints, whom they burn, murder, and destroy, yet Christ lives forever; the redemption of Christ flourishes. As a very pious poet has sung: Christ lives still, and always will; his truth also will remain, while all that fulfills this world shall perish and be worthless.

Kings perish, kingdoms perish or are changed; but the truth is never changed, Christ perishes never. He adds a most strong reason: for he is Lord of Lords, and king of kings. Therefore shall they be made a footstool for the feet of the Lamb, as many as shall strive against him. You see again why Saint John said before: they receive power as kings. For all kings are under Christ, who excels all lords in the world. For to him is given power in Heaven and in earth. Let us therefore be of bold courage. For the Lord is Emperor, and our king almighty, immortal, and invincible. He will come shortly in the clouds of fire, to judge the quick and the dead, \&c.

Moreover, victory is assuredly promised to us who are servants of Christ. And those who are with him or on his side are called, chosen, and faithful. We are chosen in Christ before the foundations of the world were laid, so that we should believe in him and be saved, first, to the Ephesians. To this we are called by the preaching of the gospel. Read 2 to the Thessalonians, chapter 2. And we ought to give thanks to God forever, and so forth. Let us hold fast to these things and be steadfast in the troubles of this world, and without fear. To God be glory.
